\section{Reinforcement Learning Game}

Creating a reinforcement learning (RL) game for a password-guessing
agent involves defining the environment, states, actions, rewards, and
possibly the policy model that will learn from interactions with the
environment. Here's a mathematical formulation of such a game:

\textbf{Environment}:

\begin{itemize}
\item
  The environment \(E\) represents the set of all possible password
  hashes \(H\) and the underlying true password distribution \(P\) from
  which these hashes are generated.
\end{itemize}

\textbf{States}:

\begin{itemize}
\item
  A state \(s \in S\) encapsulates the current knowledge of the agent
  about the distribution of passwords and any partial successes in
  guessing the passwords. It could also include a history of attempts
  and their outcomes.
\end{itemize}

\textbf{Actions}:

\begin{itemize}
\item
  The action space \(A\) includes all possible password guesses that the
  agent can make. This can be partitioned into several subsets:

  \begin{itemize}
  \item
    \(A_D\) for dictionary attack actions
  \item
    \(A_M\) for masking attack actions
  \item
    \(A_G\) for actions based on passwords generated by a GAN model.
  \end{itemize}
\end{itemize}

\textbf{Rewards}:

\begin{itemize}
\item
  The reward function \(R: S \times A \rightarrow \mathbb{R}\) provides
  immediate feedback to the agent for each action. For example,
  correctly guessing a password could yield a positive reward \(r_+\),
  while incorrect guesses could yield a negative reward \(r_-\) to
  reflect the cost of time or resources.
\end{itemize}

\textbf{Policy}:

\begin{itemize}
\item
  The policy \(\pi: S \rightarrow A\) is a mapping from states to
  actions, potentially stochastic, which the agent learns in order to
  maximize its expected cumulative reward.
\end{itemize}

\textbf{Objective}:

\begin{itemize}
\item
  The goal of the RL agent is to learn a policy \(\pi^{\ast}\) that
  maximizes the expected return from any given state over a sequence of
  \(T\) time steps, often discounted by a factor \(\gamma\). The return
  \(G_t\) from a state \(s_t\) is defined as the sum of discounted
  future rewards:
\end{itemize}

\[G_t = \sum_{k=0}^{T-t} \gamma^k R(s_{t+k}, a_{t+k})\]

\textbf{Learning Problem}:

\begin{itemize}
\item
  To find the optimal policy \(\pi^{\ast}\), we seek to maximize the
  expected return from the start state \(s_0\):
\end{itemize}

\[\pi^{\ast} = \arg\max_{\pi} \mathbb{E}_{\pi}[G_0 | s_0]\]

This expectation is taken over the distribution of states visited and
actions taken under policy \(\pi\).

\textbf{Value Functions}:

\begin{itemize}
\item
  We define the state-value function \(V^\pi(s)\) as the expected return
  starting from state \(s\) and following policy \(\pi\):
\end{itemize}

\[V^\pi(s) = \mathbb{E}_\pi[G_t | s_t = s]\]

\begin{itemize}
\item
  Similarly, the action-value function \(Q^\pi(s, a)\) is the expected
  return starting from state \(s\), taking action \(a\), and thereafter
  following policy \(\pi\):
\end{itemize}

\[Q^\pi(s, a) = \mathbb{E}_\pi[G_t | s_t = s, a_t = a]\]

The agent can utilize methods such as Q-learning, policy gradients, or
actor-critic methods to learn \(V^\pi\) or \(Q^\pi\) and improve its
policy \(\pi\).

\textbf{Learning Algorithm}:

\begin{itemize}
\item
  For instance, in Q-learning, the agent would update its estimate of
  \(Q\) using the Bellman equation as:
\end{itemize}

\[Q(s_t, a_t) \leftarrow Q(s_t, a_t) + \alpha [R(s_{t+1}, a_{t+1}) + \gamma \max_{a'}Q(s_{t+1}, a') - Q(s_t, a_t)]\]

where \(\alpha\) is the learning rate.

This RL framework can then be used to simulate the process of password
guessing, allowing the agent to learn and adapt its strategy to improve
its effectiveness over time.

\subsection{The Game}

Actions:

\begin{itemize}
    \item The type of structure to choose
\end{itemize}